\section{Dataset Sources and Construction Records}
\label{app:dataset-details}
% DOCX paragraphs 119--119, table 1. All counts are copied from the source.
\begin{table*}[t]
  \caption{GeoToolEval task sources and reference tool-invocation chain lengths. Counts refer to evaluation records.}
  \label{tab:benchmark}
  \centering
  \begin{tabularx}{\textwidth}{@{}p{0.25\textwidth}Xrrrr@{}}
    \toprule
    Task theme & Main data or question sources & 1 step & 2--3 & $\geq4$ & Total \\
    \midrule
    Three-dimensional Earth structure & LITHO1.0, SEMUCB-WM1, USGS earthquake catalog & 85 & 15 & 0 & 100 \\
    Geoscientific big-data analysis & Selected GeoQuestions1089 subset & 40 & 0 & 0 & 40 \\
    Critical minerals & Habitable-Earth questions; MRDS and cross-theme climate data & 44 & 14 & 0 & 58 \\
    Climate change & Habitable-Earth questions; paleoclimate, GISTEMP, SSP, and cross-theme mineral data & 31 & 16 & 1 & 48 \\
    Energy resources & MRDS, GEM, OPSD, World Bank, and population data & 100 & 0 & 0 & 100 \\
    Global ocean carbon cycle & SOCAT, GLODAP, SNAPO-CO$_2$ & 100 & 0 & 0 & 100 \\
    Deep-time biosphere & BioDeepTime & 12 & 7 & 5 & 24 \\
    Carbon capture and storage & Synthetic SCCS-OFPOSDU CCS--MRV data & 12 & 5 & 7 & 24 \\
    Remote sensing & Selected Earth-Bench subset & 0 & 7 & 42 & 49 \\
    Natural-disaster event statistics & EM-DAT International Disaster Database & 0 & 55 & 5 & 60 \\
    \midrule
    \textbf{Total} & & \textbf{424} & \textbf{119} & \textbf{60} & \textbf{603} \\
    \bottomrule
  \end{tabularx}
\end{table*}

\subsection{Source-Specific Task Construction}
Geoscientific big-data analysis includes 40 GeoQuestions1089 questions concerning physical geographic objects and spatial operations. The remote-sensing theme retains 49 tasks from Earth-Bench's 248-task collection, including source questions, choices, and answers, with calls adapted to the registered interfaces. Critical-mineral and climate-change questions are selected by theme from \nolinkurl{habitable_earth_dataset}; the underlying resources include MRDS, paleoclimate reconstructions, GISTEMP, and SSP warming summaries.

Earth-structure tasks use two models and an earthquake catalog. The models are LITHO1.0 and SEMUCB-WM1; earthquake records come from USGS. Energy-resource tasks use mineral, power-facility, installed-capacity, and population data. Global ocean-carbon tasks draw on SOCAT, GLODAP, and SNAPO-CO$_2$. Deep-time biosphere tasks use BioDeepTime, carbon-storage tasks use synthetic SCCS-OFPOSDU monitoring data, and disaster-event tasks use EM-DAT records. Synthetic carbon-storage data support reproducible operations under their supplied assumptions; they do not constitute field validation of storage safety.

\subsection{Evaluation Records and Provenance}
The dataset's task representation preserves one worksheet row per evaluation record. Fields identify the question and category, question text, expected answer, reference-chain information, answer type, input resource, and answer-format requirements. Stable record identities distinguish evaluation entries even when a question appears in two themes.

A source task retains its question identifier or resource locator and its reference answer. A data-derived task retains data-location information and calculations describing how the reference was obtained. Reference calls map these operational requirements to registered tools, including parameters and the relevant input/output relationships.

Construction checks record completeness and correspondence between reference-call names and registered interfaces. Sample executions check data access and invocation behavior, while computational records support recalculation. These checks are recorded separately from scored agent runs, preventing a successful construction check from being counted as an agent result.

\subsection{Interface Granularity and Aggregation}
The reference sequence is defined under the tool packaging available to the evaluated agents. A task wrapper may hide several operations inside one call, while another interface exposes them separately. Consequently, invocation count measures a property of the exposed workflow rather than an intrinsic scientific complexity.

Record-weighted averaging gives each of the 603 entries equal weight. A domain-balanced mean would instead give the ten themes equal weight, and a deduplicated question mean would change the treatment of repeated cross-theme questions. Those quantities answer different questions and are not interchangeable with the reported main FAS.

Numerical tolerances and answer normalization belong to the answer-type-specific scorer. Process metrics do not validate those numerical decisions. Conversely, matching a reference answer does not establish that the agent produced every requested artifact or followed the intended data dependencies.

The critical-mineral and climate-change themes share nine questions, represented in both themes. The 603 records therefore do not denote 603 distinct question texts. Remote sensing accounts for 42 of the 60 records with at least four reference calls, so domain and chain-length effects overlap in this dataset.

\subsection{Metric Definitions and Aggregation}
\label{app:metric-details}
Let $G_i$ be the reference sequence of tool names and $A_i$ the sequence of successful calls in the current user turn. Calls match one to one by tool name, retaining multiplicity, and $m_i$ is the number of matched occurrences. Failed calls stay out of $A_i$. Planning calls are excluded from both successful calls and attempts. A score whose denominator is zero is defined to be zero. For any per-record score $M$ and $N$ records,
\begin{equation}
\overline{M}=\frac{1}{N}\sum_{i=1}^{N}M_i.
\label{eq:metric-mean}
\end{equation}
The main experiment uses $N=603$. Domain and chain-length groups average their own records. This record-weighted mean is distinct from a domain-balanced macro-average and from the rate of fully correct answers.

FAS follows the answer type: choice letters, calendar dates, identifiers, numerical tolerances, item matching, or field comparison. An empty answer scores zero.

Tool selection is the F1 of successful-call precision and reference-call recall:
\begin{equation}
\mathrm{TSF1}_i=\frac{2m_i}{|G_i|+|A_i|}.
\label{eq:tsf1}
\end{equation}
Missing tools, extra successful calls, and repeated successful calls change the score.

Let $\pi_i$ be a longest-common-subsequence alignment from reference steps to calls in $A_i$. For reference step $k$, the argument similarity $s_{i,k}$ is zero when $k$ is not aligned. Otherwise $s_{i,k}=\mathrm{Sim}(e_{i,k},a_{i,\pi_i(k)})$, where $e_{i,k}$ is the reference argument and $a_{i,\pi_i(k)}$ is the argument of the aligned call. Equal normalized values score 1, as do scalars within a numeric tolerance. Mappings and sequences score the F1 of their recursively matched fields or items. A reference argument that is an intermediate result is compared with the value produced by the preceding successful call that supplies it. Parameter alignment is
\begin{equation}
\mathrm{PAS}_i=\frac{1}{|G_i|}\sum_{k=1}^{|G_i|}s_{i,k}.
\label{eq:pas}
\end{equation}

Tool order is the normalized longest common subsequence:
\begin{equation}
\mathrm{TOS}_i=\frac{\mathrm{LCS}(G_i,A_i)}{|G_i|}.
\label{eq:tos}
\end{equation}
For the reference order read--filter--count and the executed order read--count--filter, TSF1 is 1 and TOS is $2/3$. Extra calls do not necessarily lower TOS.

Let $U_i$ be the attempt sequence, including failed calls and excluding planning calls. Execution efficiency is the fraction of those attempts that are matched successful calls:
\begin{equation}
\mathrm{ECR}_i=\frac{m_i}{|U_i|}.
\label{eq:ecr}
\end{equation}

Quality applies fixed weights, with the largest weight on the answer and the smallest on order and attempt efficiency:
\begin{equation}
\begin{aligned}
\mathrm{Quality}_i={}&0.4\,\mathrm{FAS}_i+0.2\,\mathrm{TSF1}_i+0.2\,\mathrm{PAS}_i\\
&+0.1\,\mathrm{TOS}_i+0.1\,\mathrm{ECR}_i.
\end{aligned}
\label{eq:quality}
\end{equation}

A reference chain records one solution under the available interfaces. An alternative valid implementation can receive lower tool-name scores. PAS compares arguments only on steps aligned by tool name, so an unaligned call contributes no argument credit. These scores measure agreement with the reference workflow rather than semantic equivalence of arbitrary analyses. Explaining an answer or judging an alternative workflow still requires the inputs, outputs, and dependency trace.
